\documentclass{training}
% Filigrane PAR MODULE : \uxfond (PNG A4 dilue 30 %). Cadre noir vectoriel commun.
% Proportions original 744 px : marge 20/744, trait 1/744. Exterieur BLANC.
% Cadre via eso-pic +\put (coordonnees absolues) — pas de remember picture (decalage).
\usepackage{eso-pic}

\providecommand{\uxfond}{../../resources/fond/norme_fond_paliers_30.png}

\newlength{\uxmargin}
\newlength{\uxthick}
\newlength{\uxinnerw}
\newlength{\uxinnerh}
\setlength{\uxmargin}{0.02688\paperwidth}% 20/744
\setlength{\uxthick}{0.001344\paperwidth}% 1/744
\setlength{\uxinnerw}{\dimexpr\paperwidth-2\uxmargin-2\uxthick\relax}
\setlength{\uxinnerh}{\dimexpr\paperheight-2\uxmargin-2\uxthick\relax}

% Contenu strictement a l'interieur du trait (+ jour confortable).
\geometry{%
  includeheadfoot,
  left=22mm,
  right=22mm,
  top=20mm,
  bottom=18mm,
  headheight=30pt,
  headsep=5mm,
  footskip=12mm%
}

\AddToShipoutPictureBG{%
  \AtPageLowerLeft{%
    % Fond blanc page entiere
    \put(0,0){\color{white}\rule{\paperwidth}{\paperheight}}%
    % Filigrane module a l'interieur du cadre
    \put(\LenToUnit{\dimexpr\uxmargin+\uxthick\relax},\LenToUnit{\dimexpr\uxmargin+\uxthick\relax}){%
      \includegraphics[width=\uxinnerw,height=\uxinnerh]{\uxfond}}%
    % Cadre noir : bas, haut, gauche, droite
    \put(\LenToUnit{\uxmargin},\LenToUnit{\uxmargin}){%
      \color{black}\rule{\dimexpr\paperwidth-2\uxmargin\relax}{\uxthick}}%
    \put(\LenToUnit{\uxmargin},\LenToUnit{\dimexpr\paperheight-\uxmargin-\uxthick\relax}){%
      \color{black}\rule{\dimexpr\paperwidth-2\uxmargin\relax}{\uxthick}}%
    \put(\LenToUnit{\uxmargin},\LenToUnit{\uxmargin}){%
      \color{black}\rule{\uxthick}{\dimexpr\paperheight-2\uxmargin\relax}}%
    \put(\LenToUnit{\dimexpr\paperwidth-\uxmargin-\uxthick\relax},\LenToUnit{\uxmargin}){%
      \color{black}\rule{\uxthick}{\dimexpr\paperheight-2\uxmargin\relax}}%
  }%
}

\makeatletter
\fancyhf{}
\fancyheadoffset{0pt}
\renewcommand{\headrulewidth}{0pt}
\lhead{\parbox[t]{0.56\textwidth}{\RaggedRight\footnotesize\sloppy\@title}}
\rhead{\parbox[t]{0.36\textwidth}{\RaggedLeft\footnotesize\sloppy\@subtitle}}
\cfoot{\thepage}
% Coupures : underscore et slash dans les noms de fichiers / consignes
\DeclareRobustCommand{\_}{\textunderscore\allowbreak}
\makeatother

% Evite les depassements de longues lignes (Github, Name_FirstName_...)
\sloppy
\emergencystretch=3em
\hbadness=10000

% Sommaire type 42 : chapitres seulement
\setcounter{tocdepth}{0}
\renewcommand{\cftchappresnum}{}
\renewcommand{\cftchapaftersnum}{}
\setlength{\cftchapnumwidth}{2.0em}
\setlength{\cftbeforechapskip}{0.35em}
\cftsetindents{chapter}{0em}{2.0em}


\begin{document}

\title{Seamstress}
\subtitle{Piscine UX/UI --- UX}
\member{Piscine UX/UI}{https://github.com/01-edu/UX-UI}
\summary{Strategy, empathy, definition, statement, ideas, prototype, test.}
\version{1.0}
\prerequisite{The previous module of this piscine, or none if this is your first.}
\maketitle

\newpage
\tableofcontents
\newpage


\chapter{Instructions}

The foreword is not an exercise. It will not be graded. The exercises start at the next chapter, from the easiest contract to the hardest. A later exercise is not taken into account if an earlier one does not hold.

Only what is in the turn-in directory of an exercise is looked at. Do not add files that the exercise does not name. The name of a deliverable follows the pattern in the exercise: \texttt{Name\_FirstName\_DeliverableName\_Date\_VersionNumber}.

You will not be given the gesture. The examples in an exercise sometimes carry a constraint that the sentence does not repeat. Read them. Your other references are the manual of the tool you picked, the person sitting next to you, and the links printed in the exercise.

The order is the order of the quest. Skipping a step because a later screen looks more interesting produces a screen nobody can explain.

\subsection*{About the pictures}

The picture in the foreword is not a hint. The pale steps behind the text are the page ground.

\subsection*{About machines that write}

You may use a tool to look up a term, the way you would use a dictionary. You may not replace an interview, a test, or a file with a generated text and sign it. If you cannot point at the person, the screen, or the line you claim, the exercise is not done.


\chapter{Foreword}

The Design Council published the double diamond as a map: two widenings, two narrowings (\url{https://www.designcouncil.org.uk/our-resources/the-double-diamond/}). Charles and Ray Eames liked to say that the details are not the details, they make the product. This quest walks the diamond. The map is not the garment.

\begin{center}
\includegraphics[width=0.78\textwidth,height=0.32\textheight,keepaspectratio]{../../resources/jeunecrack.png}\\[0.4em]
{\footnotesize Picture: JeuneCrack}
\end{center}

Recommended watching. The videos are context. They are not the deliverable.
\begin{itemize}
\item UX - Introduction \& Quest n°1, 01Edu \url{https://www.youtube.com/watch?v=UbE6WMeuvQg&list=PLHyAJ_GrRtf-RAtfngbRrFh1XDTnyXFDe}
\end{itemize}

\chapter{Exercise 00: UX VII - Ex 1 - UX Strategy}
\exnumber{00}
\extitle{UX VII - Ex 1 - UX Strategy}
\exfiles{ex00/Figma file and PDF, named Name\_FirstName\_DeliverableName\_Date\_VersionNumber}
\exauthorize{Those named in the exercise}
\makeheaderfiles

\textbf{Context:}

Elsa is a seamstress. She makes swimsuits on demand through her instagram account. It's going well, her community is growing as well as the amount of requests from customers.

However, she loses a lot of time every time there is a request. To complete an order, she needs to ask for the shape of the swimsuit, the size, the color preferences and the expected delivery time. On top of that, she needs to make sure she has enough time to sew one piece - every piece requires 6 hours.

For now, she only works with private messages but she wishes she could spend less time on that part.

You offer to help her design a solution to manage her time and workshop better. You have 4 working days.

\textbf{Instructions:}

Now you have discovered and practiced many tools in the design process, organize a time frame for the next 4 working days.

Select at least one tool for each phase, name it, and take into account the amount of time each phase may require.

\begin{itemize}
\item Empathy / User Research
\item Define
\item Problem Statement
\item Ideation
\item Prototype
\item Test
\end{itemize}

\textbf{Deliverables:}

\begin{itemize}
\item A time frame, on PDF or on Figma, of the project, including the working days ahead, each phase of the process, the main tools you plan to use for each phase and the time you plan to spend for each tool.
\item Upload the PDF or Figma document on Github, labeled as follows: "Name\_FirstName\_DeliverableName\_Date\_VersionNumber". For example, it can be labeled as "Doe\_John\_MidFiPrototype\_05242024\_V1".
\item Add a title within the document.
\item Pay attention to the global aspect of the document. It must be clear, simple and easy to read. You can get inspiration on canva.com but don't overload your design with too much details!
\end{itemize}

\textbf{Resource:}

\begin{itemize}
\item A sum up of all the tools that have been practiced \url{https://www.figma.com/file/x0xfoz9TaEDiAfGPFqKJcs/RECAP-UX%2FUI-Program?node-id=0%3A1}. Focus on the UX part!
\end{itemize}

\textbf{Quotes:}

\begin{itemize}
\item Design needs to fail. Failure is even a necessary step, but ideally it should happen before a product is launched, during the prototype and test phases
\item "UX Strategy lies at the intersection of UX and business. It provides a much better chance of creating successful products. It enables teams to see the "Big Picture" to achieve the business goals under uncertain conditions" Jamie Levy | UX Strategy: How to Devise Innovative Digital Products That People Want
\item UX Strategy is the method by which you validate that your solution solves a problem for real customers in a dynamic marketplace because the market is constantly changing.
\end{itemize}


\chapter{Exercise 01: UX VII - Ex 2 - Empathy}
\exnumber{01}
\extitle{UX VII - Ex 2 - Empathy}
\exfiles{ex01/zip archive, named Name\_FirstName\_DeliverableName\_Date\_VersionNumber}
\exauthorize{Those named in the exercise}
\makeheaderfiles

\textbf{Instructions:}

Run at least one of the Empathy tools in the design process.

\textbf{Deliverables:}

\begin{itemize}
\item A PDF document with the script of the tool context
\item Completed tool (scripts from user interviews, results from a quantitative analysis, etc.)
\item A PDF document with the key learnings from this phase
\item Upload the documents on Github, in a zip folder named "ProjectTitle\_Name\_FirstName", with all project deliverables as follows: "Name\_FirstName\_DeliverableName\_Date\_VersionNumber".  For example, the first deliverable can be named as "Doe\_John\_MidFiPrototype\_05242024\_V1".
\item Add a title within all written documents.
\item Pay attention to the global aspect of the deliverables. They must be clear, simple and easy to read. You can get inspiration on canva.com but don't overload your design with too much details!
\end{itemize}

\textbf{Tips:}

\begin{itemize}
\item Use as many tools as needed, be make sure you do all the process in 4 working days!
\end{itemize}

\textbf{Criteria:}

\begin{itemize}
\item User interviews
\begin{itemize}
\item The script is the interview guide. Make sure there is:
\begin{itemize}
\item the main purpose of the interview
\item 1 to 3 hypotheses
\item a proper introduction that describes the interviewer and the purpose of this interview
\item 3 to 5 main topics
\item a few scan questions (demographics)
\item open-ended questions, max 10
\item a funnel process (>)
\item a conclusion with next steps and a thank you.
\end{itemize}
\item The completed tool is either:
\begin{itemize}
\item 5 to 10 scripts from users
\item or an Affinity Diagram \url{https://www.youtube.com/watch?v=VngWHIE4k9s} with the main elements
\end{itemize}
\end{itemize}
\end{itemize}

\begin{itemize}
\item Analytics
\begin{itemize}
\item The script is the form. Make sure there is:
\begin{itemize}
\item The main purpose of the form
\item 1 to 3 hypotheses
\item A proper introduction that describes the interviewer and the purpose of this interview
\item 3 to 5 main topics
\item A few scan questions (demographics)
\item Closed-ended, unbiased questions, max 10
\item A funnel process (>)
\item A conclusion with next steps and a thank you.
\item A broadcast strategy
\end{itemize}
\item The completed tool is a series of quantitative data from at least 100 answers.
\end{itemize}
\end{itemize}

\begin{itemize}
\item Competitive analysis
\begin{itemize}
\item The script is the strategy. The student explains what they are looking for and how they plan to find it.
\item The completed tool is a report. Make sure there is;
\begin{itemize}
\item There is an introduction
\item There is a competitive landscape \url{https://venngage-wordpress.s3.amazonaws.com/uploads/2019/03/g2crowd-grid.png}
\item With 2 relevant axes
\item 10 organizations appear on the landscape
\item There is a competitive table \url{https://venngage-wordpress.s3.amazonaws.com/uploads/2019/03/product-features-comparison-matrix.png}
\item With 5 to 10 criteria
\item 10 organizations appear on the table
\end{itemize}
\end{itemize}
\end{itemize}


\chapter{Exercise 02: UX VII - Ex 3 - Define}
\exnumber{02}
\extitle{UX VII - Ex 3 - Define}
\exfiles{ex02/zip archive, named Name\_FirstName\_DeliverableName\_Date\_VersionNumber}
\exauthorize{Those named in the exercise}
\makeheaderfiles

\textbf{Instructions:}

Run at least one of the Define tools in the design process.

\textbf{Deliverables:}

\begin{itemize}
\item A PDF document with the script of the tool context
\item Completed tool (scripts from user interviews, results from a quantitative analysis, etc.)
\item A PDF document with the key learnings from this phase
\item Upload the documents on Github, in a zip folder named "ProjectTitle\_Name\_FirstName", with all project deliverables as follows: "Name\_FirstName\_DeliverableName\_Date\_VersionNumber".  For example, the first deliverable can be named as "Doe\_John\_MidFiPrototype\_05242024\_V1".
\item Add a title within all written documents.
\item Pay attention to the global aspect of the deliverables. They must be clear, simple and easy to read. You can get inspiration on canva.com but don't overload your design with too much details!
\end{itemize}

\textbf{Tips:}

\begin{itemize}
\item Use as many tools as needed, be make sure you do all the process in 4 working days!
\end{itemize}

\textbf{Criteria:}

\begin{itemize}
\item Persona
\begin{itemize}
\item The script is the list of insights the students uses to create the persona
\item The completed tool is a Persona board \url{https://miro.medium.com/max/1400/0*RbHfR6BHhhFIENNf}. Make sure the key insights come from the field, from reality, but the persona itself is different from only one user.
\end{itemize}
\end{itemize}

\begin{itemize}
\item Empathy Map
\begin{itemize}
\item The script is the list of insights the student uses to create the Empathy map
\item The completed tool is an Empathy Map \url{https://miro.medium.com/max/700/1*I1ffOWdPWQva3dCMQE-TAQ.png}.
\end{itemize}
\end{itemize}

\begin{itemize}
\item User Journey
\begin{itemize}
\item The script is the list of insights the student uses to create the User Journey.
\item The completed tool is a User Journey Map \url{https://d2slcw3kip6qmk.cloudfront.net/marketing/blog/2017Q3/SEO-initiative-customer-journey-mapping/CustomerJourneyMap1.png}.
\end{itemize}
\end{itemize}

\begin{itemize}
\item Jobs-to-be-done (JTBD)
\begin{itemize}
\item The JTBD is tool that goes as a conclusion to another of the previous tools. Make sure:
\begin{itemize}
\item The sentence is complete
\item The "When" part involves a visit on the website
\item The "I want to" part involves a practical action
\item The "So I can" part involves a psychological or emotional purpose
\end{itemize}
\end{itemize}
\end{itemize}


\chapter{Exercise 03: UX VII - Ex 4 - Problem Statement}
\exnumber{03}
\extitle{UX VII - Ex 4 - Problem Statement}
\exfiles{ex03/PDF, named Name\_FirstName\_DeliverableName\_Date\_VersionNumber}
\exauthorize{Those named in the exercise}
\makeheaderfiles

\textbf{Instructions:}

Based on your research (Empathy \& Define), define a problem statement to narrow the problem and KPIs to measure the success.

\textbf{Deliverables:}

\begin{itemize}
\item A PDF document with:
\begin{itemize}
\item One problem statement
\item 1 to 3 KPIs
\item Key learnings from this phase
\end{itemize}
\item Upload the PDF document on Github, labeled as follows: "Name\_FirstName\_DeliverableName\_Date\_VersionNumber". For example, it can be labeled as "Doe\_John\_MidFiPrototype\_05242024\_V1".
\item Add a title within the document.
\item Pay attention to the global aspect of the document. It must be clear, simple and easy to read. You can get inspiration on canva.com but don't overload your design with too much details!
\end{itemize}

\textbf{Tips:}
\begin{itemize}
\item KPIs: check the SMART list \url{https://www.grow.com/blog/how-to-use-smart-goals-to-build-your-kpis}) 
\end{itemize}

\textbf{Quote:}

Design is about solving problems. Fall in Love with the Problem, Not the Solution.

\textbf{Examples of Problem statements:}

\begin{itemize}
\item How might we help Pierre measure and prepare his hydration before a race?
\begin{itemize}
\item Appropriate  --- > The problem statement is narrow enough to be specific for one painpoint, and broad enough to let creativity suggest many options. Also, it takes into accounts some of the key insights from the previous phases.
\end{itemize}
\item How might we make sure Pierre doesn't get dehydrated?
\begin{itemize}
\item Too broad  --- > The problem statement is out of the scope of running, so the solutions cannot be integrated in the specific context of running.
\end{itemize}
\item How might we make sure Pierre drinks the cups that are offered during the race?
\begin{itemize}
\item Too narrow  --- > There are plenty of options to avoid dehydration without drinking the cups that are offered during the race.
\end{itemize}
\item How might we provide Pierre with a drink that contains everything he needs to avoid dehydration?
\begin{itemize}
\item Focused on the solution, not the problem  --- > The problem statement is supposed to rephrase the brief and take into account the insights from the previous phase.
\end{itemize}
\item How might make sure Pierre doesn't get drunk the night before the race?
\begin{itemize}
\item Out of scope  --- > This is not about dehydration, nor about running.
\end{itemize}
\end{itemize}


\chapter{Exercise 04: UX VII - Ex 5 - Ideation}
\exnumber{04}
\extitle{UX VII - Ex 5 - Ideation}
\exfiles{ex04/Figma file and PDF, named Name\_FirstName\_DeliverableName\_Date\_VersionNumber}
\exauthorize{Those named in the exercise}
\makeheaderfiles

\textbf{Instructions:}

Run at least one of the Ideation techniques in the design process.

\textbf{Deliverables:}

\begin{itemize}
\item A PDF document with the script of the tool context
\item Completed tool: a board with all the ideas - on paper or digitally on Miro or Figma
\item A PDF document with the key learnings from this phase
\item Upload the documents on Github, in a zip folder named "ProjectTitle\_Name\_FirstName", with all project deliverables as follows: "Name\_FirstName\_DeliverableName\_Date\_VersionNumber".  For example, the first deliverable can be named as "Doe\_John\_MidFiPrototype\_05242024\_V1".
\item Add a title within all written documents.
\item Pay attention to the global aspect of the deliverables. They must be clear, simple and easy to read. You can get inspiration on canva.com but don't overload your design with too much details!
\end{itemize}

\textbf{Tips:}

\begin{itemize}
\item Use as many tools as needed, and make sure you do all the process in 4 working days!
\end{itemize}


\chapter{Exercise 05: UX VII - Ex 6 - Prototype}
\exnumber{05}
\extitle{UX VII - Ex 6 - Prototype}
\exfiles{ex05/zip archive, named Name\_FirstName\_DeliverableName\_Date\_VersionNumber}
\exauthorize{Those named in the exercise}
\makeheaderfiles

\textbf{Instructions:}

Make a prototype including the chosen solution.

\textbf{Deliverables:}

\begin{itemize}
\item Pictures of a paper prototype (low-fi)
\item An animated wireframe prototype (mid-fi)
\item Key learnings from this phase
\item Vocabulary
\begin{itemize}
\item Low-Fidelity: Paper Prototypes - Example here \url{https://miro.medium.com/max/1400/1*5nUfqqA2gjdAYHagjbHA5w.jpeg}
\item Mid-Fidelity: Wireframes - black and white, only the main elements, on digital prototype tools - Example here \url{https://miro.medium.com/max/1400/1*eGSfaxwYbxUFomYbyO6GWw.png}
\item High-Fidelity: Wireframes with color, styles, graphical details, and micro-interactions - Example here \url{https://miro.medium.com/max/1400/1*Xn0HSKAvhr4TZzC9lN5udw.gif}
\end{itemize}
\item Upload the documents on Github, in a zip folder named "ProjectTitle\_Name\_FirstName", with all project deliverables as follows: "Name\_FirstName\_DeliverableName\_Date\_VersionNumber".  For example, the first deliverable can be named as "Doe\_John\_MidFiPrototype\_05242024\_V1".
\item Add a title within all written documents.
\item Pay attention to the global aspect of the deliverables. They must be clear, simple and easy to read. You can get inspiration on canva.com but don't overload your design with too much details!
\end{itemize}

\textbf{Tips:}

\begin{itemize}
\item Use as many tools as needed, and make sure you do all the process in 4 working days!
\item You may use Card Sorting and User flows
\end{itemize}

\textbf{Criteria:}

\subsubsection*{Criteria}

\begin{itemize}
\item User Flow
\end{itemize}

Make sure for each flow:

\begin{itemize}
\item The user has a specific goal
\item There is a starting point and and ending point - which is the goal of the user
\item There is a series of actions organised in rectangles (actions) or diamond (choice)
\item The actions are linked with arrows
\end{itemize}

\begin{itemize}
\item Card sorting
\end{itemize}

Make sure

\begin{itemize}
\item There is a script including the research goals.
\item The research goals must be close to "rearrange the content"
\item There are pictures of the ongoing workshop
\item There is a new content organization shown in a table with sticky notes (it can be a picture or digital)
\end{itemize}


\chapter{Exercise 06: UX VII - Ex 7 - Test}
\exnumber{06}
\extitle{UX VII - Ex 7 - Test}
\exfiles{ex06/zip archive, named Name\_FirstName\_DeliverableName\_Date\_VersionNumber}
\exauthorize{Those named in the exercise}
\makeheaderfiles

\textbf{Instructions:}

Prepare a testing protocol and run it!

\textbf{Deliverables:}

\begin{itemize}
\item A testing protocol script
\item Photos from the tests
\item Key learnings from the tests (5 tests are required)
\item Key learnings from this phase
\item Upload the documents on Github, in a zip folder named "ProjectTitle\_Name\_FirstName", with all project deliverables as follows: "Name\_FirstName\_DeliverableName\_Date\_VersionNumber".  For example, the first deliverable can be named as "Doe\_John\_MidFiPrototype\_05242024\_V1".
\item Add a title within all written documents.
\item Pay attention to the global aspect of the deliverables. They must be clear, simple and easy to read. You can get inspiration on canva.com but don't overload your design with too much details!
\end{itemize}


\end{document}
